\documentclass[11pt]{article}
\usepackage{amsmath,amssymb,amsthm}

\newtheorem{theorem}{Theorem}
\newtheorem{corollary}[theorem]{Corollary}
\newtheorem{lemma}[theorem]{Lemma}

\newcommand{\Up}{U_+}
\newcommand{\mass}{\operatorname{B}}
\newcommand{\car}{\operatorname{C}}
\newcommand{\sh}{\operatorname{D}}

\title{Arbitrary-prefix same-depth inverse rigidity on a record-atom tail}
\author{}
\date{27 August 2026}

\begin{document}
\maketitle

For a positive odd integer $x$, put
\[
 \Up(x):=\frac{3x+1}{2^{v_2(3x+1)}}.
 \tag{1}
\]
For an exponent word $w=(a_1,\ldots,a_L)$, define
\[
 \mass(w):=\sum_{i=1}^L a_i,
 \qquad
 \car(w):=\sum_{i=1}^L3^{L-i}2^{a_1+\cdots+a_{i-1}}.
 \tag{2}
\]
Thus an exact path with word $w$ satisfies
\[
 2^{\mass(w)}\Up^L(x)=3^Lx+\car(w).
 \tag{3}
\]
Use the record-atom notation
\[
 A:=1,\qquad B:=12,\qquad C:=1113,\qquad D:=1122.
 \tag{4}
\]
The purpose of this note is to classify the least same-depth inverse over
an arbitrary finite word in these atoms.  No bound on the number of atoms
or on the intervening runs of $A$ is imposed.

\section{Shifted carries}

It is convenient to translate every state by one.  Define
\[
 \sh(w):=\car(w)+2^{\mass(w)}-3^{|w|}.
 \tag{5}
\]
If $X=x+1$ and $Y=\Up^{|w|}(x)+1$, then (3) becomes
\[
 2^{\mass(w)}Y=3^{|w|}X+\sh(w).
 \tag{6}
\]

\begin{lemma}[nonnegative and uniformly bounded shifted carry]
For every positive exponent word $v=(d_1,\ldots,d_\ell)$,
\[
 \sh(v)=\sum_{i=1}^{\ell}
 3^{\ell-i}2^{d_1+\cdots+d_{i-1}}(2^{d_i}-2)\geq0
 \tag{7}
\]
and
\[
 \frac{\car(v)}{2^{\mass(v)}}
 \leq\sum_{j=0}^{\ell-1}\frac{3^j}{2^{j+1}}
 =\left(\frac32\right)^\ell-1.
 \tag{8}
\]
If $W$ is a nonempty finite word over the atoms $A,B,C,D$, then
\[
 \boxed{0\leq\sh(W)<4\,3^{|W|}.}
 \tag{9}
\]
\end{lemma}

\begin{proof}
Expanding (5) one letter at a time gives (7).  Also
\[
 \frac{\car(v)}{2^{\mass(v)}}
 =\sum_{i=1}^{\ell}
 \frac{3^{\ell-i}}{2^{d_i+\cdots+d_\ell}},
\]
and $d_i+\cdots+d_\ell\geq\ell-i+1$, proving (8).

For concatenated words,
\[
 \sh(UV)=3^{|V|}\sh(U)+2^{\mass(U)}\sh(V).
 \tag{10}
\]
The atom data are
\[
\begin{array}{c|c|c|c|c}
P&|P|&\mass(P)&\car(P)&\sh(P)\\ \hline
A&1&1&1&0\\
B&2&3&5&4\\
C&4&6&65&48\\
D&4&6&73&56.
\end{array}
\tag{11}
\]
Write $h(W)=\sh(W)/3^{|W|}$.  From (10),
\[
 h(PU)=\frac{\sh(P)}{3^{|P|}}
       +\frac{2^{\mass(P)}}{3^{|P|}}h(U).
 \tag{12}
\]
If $h(U)<4$, the four right sides in (12) are respectively less than
\[
 \frac83,\qquad
 \frac4{9}+\frac{32}{9}=4,\qquad
 \frac{48+256}{81},\qquad
 \frac{56+256}{81}.
\]
They are all at most $4$, with strictness inherited in the $B$ case.
Induction on the number of atoms proves (9).
\end{proof}

\section{Distortion states}

Let $W$ be an exact record-atom word from $x$ to $y$, and let $V$ be
an exact positive exponent word of the same raw length from $z$ to $y$.
At any common raw-time boundary put
\[
 \Delta:=\mass(V_{\rm rem})-\mass(W_{\rm rem}),
 \qquad z=2^\Delta x+c,
 \qquad u:=\frac{c}{2^\Delta},
 \tag{13}
\]
where ``rem'' denotes the two remaining suffixes from that boundary to
the common endpoint.  At the endpoint the state is $(\Delta,c)=(0,0)$.

Suppose an actual atom $P$ of length $\ell$, mass $b$, and carry $C_P$
is prepended to a processed suffix.  Partition the alternative path at
the same raw time and let $v$ be the corresponding length-$\ell$ block,
with mass $s$ and carry $C_v$.  If $(\Delta,c)$ is the state on the right
and $(\Delta',c')$ the state on the left, then
\[
 \boxed{
 \Delta'=\Delta+s-b,
 \qquad
 c'=\frac{2^{\Delta'}C_P+2^s c-C_v}{3^\ell}.}
 \tag{14}
\]
Consequently, with $R(v):=C_v/2^s$,
\[
 \boxed{
 u'=\frac{2^b}{3^\ell}u+\frac{C_P}{3^\ell}
       -\frac{2^{b-\Delta}}{3^\ell}R(v).}
 \tag{15}
\]

Because both boundary states are odd, every genuine state lies in the
parity lattice
\[
 \begin{cases}
 c\in2\mathbb Z+1,&\Delta\geq1,\\
 c\in2\mathbb Z,&\Delta=0,\\
 c=m/2^{-\Delta}\text{ with }m\text{ odd},&\Delta<0.
 \end{cases}
 \tag{16}
\]

The only finite arithmetic needed below is isolated next.  It is included
as an exact certificate rather than hidden in a search over word length.

\newpage
\begin{lemma}[exceptional transition certificate]
For a lattice-admissible left state with $\Delta'\geq-1$, the eight
exceptional cases have the following sharp bounds; a dash means none exists.
\[
\begin{array}{c|c|c|c}
(\Delta,c)&P&\max R(v)&\text{maximizing }v\\ \hline
(-1,-\tfrac12)&B&\text{--}&\text{--}\\
(0,0)&B&5/8&12\\ \hline
(-1,-\tfrac12)&C&85/512&1215\\
(-1,-\tfrac12)&D&\text{--}&\text{--}\\
(0,0)&C&73/32&1121\\
(0,0)&D&73/64&1122\\
(1,-1)&C&\text{--}&\text{--}\\
(1,-1)&D&85/256&1214.
\end{array}
\tag{17}
\]
If the right state satisfies $\Delta\geq-1$ and $-1\leq u\leq5$, with
$c=-1/2$ when $\Delta=-1$, then no lattice-admissible transition has
$\Delta'<-1$.
Moreover, among all strip states which enter the half-state
$(-1,-1/2)$ in one atom, the complete preimage table is
\[
\begin{array}{c|c|c|c}
(\Delta,c)_{\rm right}&P&v&(\Delta',c')_{\rm left}\\ \hline
(-1,-\tfrac12)&A&1&(-1,-\tfrac12)\\
(0,0)&C&1121&(-1,-\tfrac12).
\end{array}
\tag{18}
\]
\end{lemma}

\begin{proof}
For $\Delta'\leq2$, the total mass
$s=\Delta'-\Delta+b$ is fixed.  Substitution of all positive
compositions of $s$ into (14) gives the low part of (17).

For $\Delta'\geq3$, write $v=(a,\eta)$.  In all eight cases the numerator
in (14) is an integer of the form
\[
 2^aK(\eta)-3^{\ell-1}.
 \tag{19}
\]
Thus divisibility by $3^\ell$ depends on $a$ only through its parity,
because increasing $a$ by two changes (19) from $N$ to
$4N+3^\ell$.  The suffix exponents may be reduced to positive residues
modulo
\[
 \operatorname{ord}_{9}(2)=6,
 \qquad
 \operatorname{ord}_{81}(2)=54.
 \tag{20}
\]
Indeed, transfer of one full period from a suffix exponent to $a$
preserves the total mass and the congruence, and can only increase
\[
 R(v)=\sum_{i=1}^{\ell}
 \frac{3^{\ell-i}}{2^{d_i+\cdots+d_\ell}}.
 \tag{21}
\]
\newpage
For each of the two parities, lowering $a$ to the first positive value
with $\Delta'\geq3$ again preserves the congruence and increases (21).
It therefore suffices to inspect $6$ suffix residues for $B$ and
$54^3$ suffix triples for $C,D$, with two choices of parity.  Exact
integer substitution gives (17).

The cases which could have $\Delta'<-1$ have total mass at most five;
direct substitution in (14) gives no point of the lattice (16).
Finally, imposing $(\Delta',c')=(-1,-1/2)$ bounds the right defect by
$1$ and the alternative mass by $6$.  Substitution gives exactly the two
rows of (18).  The accompanying program
\texttt{verify\_divergence\_record\_arbitrary\_prefix.py} performs these
finite integer checks and independently verifies every maximum in (17).
\end{proof}

\begin{lemma}[uniform distortion strip and rigidity at defect $-1$]
At every atom boundary of the two aligned paths,
\[
 \boxed{\Delta\geq-1,\qquad -1\leq u\leq5.}
 \tag{22}
\]
If the remaining actual suffix is nonempty and $\Delta=-1$, then
\[
 \boxed{c=-\frac12,\qquad z=\frac{x-1}{2}.}
 \tag{23}
\]
\end{lemma}

\begin{proof}
We induct from the common endpoint toward the two starts.  The empty
suffix has $(\Delta,c,u)=(0,0,0)$.

First consider the upper edge of the strip.  Since $C_v>0$, equation
(15) and $u\leq5$ give
\[
 u'<\frac{5\,2^b+C_P}{3^\ell}\leq5;
 \tag{24}
\]
the four numerators for $A,B,C,D$ are $11,45,385,393$, while the four
denominators are $3,9,81,81$.

For the lower edge, (8) gives $R(v)\leq1/2,5/4,65/16$ for blocks of
length $1,2,4$, respectively.  Substitution in (15) proves $u'\geq-1$
directly for
\[
 A:\Delta\geq-1,
 \qquad B:\Delta\geq1,
 \qquad C,D:\Delta\geq2.
 \tag{25}
\]
At the smallest listed defects the resulting lower bounds are
\[
 -1,\qquad -\frac89,\qquad -\frac{64}{81},\qquad -\frac{56}{81}.
\]

\newpage
It remains to treat lower right defects.  By the inductive form of (23),
$\Delta=-1$ has $c=-1/2$.  If $\Delta=0$, the lattice and the strip give
$c\in\{0,2,4\}$; if $\Delta=1$, they give
$c\in\{-1,1,3,5,7,9\}$.  The worst values $c=-1/2,0,-1$ are precisely
the rows of (17).  Equation (15) then gives, in the nonempty cases, the
lower bounds
\[
 0,-\frac14,-1,0,\frac38
 \tag{26}
\]
for $B(0),C(-1),C(0),D(0),D(1,-1)$.
For $c\geq2$ at defect zero and $c\geq1$ at defect one, the unrestricted
bound (8) already gives a value greater than $-1$.  The finite
below-$-1$ assertion of Lemma~2 proves $\Delta'\geq-1$.  This establishes
(22).

Suppose now that a nonempty processed suffix $W$ has $\Delta=-1$.
The lattice and (22) give
\[
 2c\in\{-1,1,3,5\}.
 \tag{27}
\]
Put $X=x+1$, $Z=z+1$, and
\[
 k:=-c-\frac12,
 \qquad Z=\frac{X}{2}-k.
 \tag{28}
\]
Since the alternative suffix has one less unit of mass, its shifted
equation (6), multiplied by two and compared with the actual equation,
gives
\[
 \sh(V)=\frac{\sh(W)}2+k3^{|W|}.
 \tag{29}
\]
By Lemma~1, $\sh(V)\geq0$ and $\sh(W)<4\,3^{|W|}$.  Therefore
$k\leq-2$ is impossible.  The value $k=-1$ corresponds to $c=1/2$,
but every atom begins with exponent one, so $x\equiv3\pmod4$; then
$z=(x+1)/2$ is even.  The only remaining value is $k=0$, which is (23).
This also closes the inductive assertion used above.
\end{proof}

\section{The arbitrary-prefix minimum}

\begin{theorem}[arbitrary-prefix same-depth inverse minimum]
Let
\[
 W=PU,
 \qquad P\in\{B,C,D\},
 \qquad U\in\{A,B,C,D\}^*,
 \tag{30}
\]
be an exact word from a positive odd integer $q$ to $y$, and put
$L=|W|$.  Among all positive odd $L$-step predecessors of $y$, the
unique least one is
\[
 \boxed{
 q_{\min}=\begin{cases}
 q,&P=B\text{ or }D,\\[1mm]
 (q-1)/2,&P=C.
 \end{cases}}
 \tag{31}
\]
In the second case the unique minimum word is $1121\,U$; it is exactly
the local half-coalescing branch at the initial atom $C$, followed by the
unchanged suffix $U$.
\end{theorem}

\begin{proof}
Let $V$ be any exact length-$L$ word from $z$ to $y$, and form its start
state $(\Delta,c)$.  Lemma~3 gives $\Delta\geq-1$.

If $\Delta=-1$, equation (23) gives $z=(q-1)/2$.  Read the first actual
atom forward, or equivalently use the one-atom preimage table (18).
Since the left state is the half-state, the first atom is either $A$ and
the half-state persists, or it is $C$, the alternative block is $1121$,
and the right state is $(0,0)$.  Here the first atom belongs to
$\{B,C,D\}$, so only the second alternative is possible.  Thus $P=C$.
At the right boundary the two positive odd states are equal.  The
accelerated map is deterministic, so the remaining exact word is $U$ on
both paths.  Conversely, exactness of $C$ gives
$q\equiv79\pmod {128}$; direct substitution then gives the exact word
$1121$ from $(q-1)/2$ and the identity
\[
 \Up^4\!\left(\frac{q-1}{2}\right)=\Up^4(q)
 \tag{32}
\]
for the words $1121$ and $1113$ shows that this branch always exists
when $P=C$.

If $\Delta=0$, the lattice and (22) give $c\in\{0,2,4\}$, hence
$z=q+c\geq q$.  Equality forces the same start and therefore the same
deterministic word $V=W$.  If $\Delta\geq1$, (22) gives
\[
 z=2^\Delta q+c\geq2^\Delta(q-1)>q,
 \tag{33}
\]
because an atom start has $q\equiv3\pmod4$ and hence $q\geq3$.
These alternatives prove (31) and uniqueness.
\end{proof}

\enlargethispage{3\baselineskip}
\begin{corollary}[no finite same-depth pullback creates a new descent]
Suppose an interior record tail contains a finite block beginning at a
nontrivial atom $P\in\{B,C,D\}$ and ending at any later atom boundary.
Extending the block across arbitrarily many intervening record atoms can
never create a new same-depth smaller predecessor.  If $P=B$ or $D$, the
actual block start remains uniquely least.  If $P=C$, the unique smaller
predecessor is only the local half-seed at that initial $C$.

Consequently, pulling a terminal-$C$ half-coalescence across any finite
number of preceding record atoms cannot yield a new same-depth descent.
Any genuinely new descent on this branch must change the raw inverse
depth.
\end{corollary}

\begin{proof}
Apply Theorem~4 to the complete finite atom word.  Its two alternatives
depend only on the first nontrivial atom, not on the length or composition
of the suffix.
\end{proof}

\end{document}
